\documentclass[12pt]{article}
\usepackage[utf8]{inputenc}
\usepackage[T1]{fontenc}
\usepackage[margin=1in]{geometry}
\usepackage{setspace}
\usepackage{xcolor}
\usepackage[colorlinks=true, urlcolor=blue!60!black, linkcolor=blue!60!black]{hyperref}
\setlength{\parindent}{0em}
\setlength{\parskip}{0.6em}
\renewcommand{\baselinestretch}{1.0}
\begin{document}
\begin{center}
{\LARGE \textit{Fitzpatrick v. Bitzer}, 427 U.S. 445 (1976)\par}
\vspace{0.5em}
{\large Case Brief --- Two-Page Edition\par}
\end{center}

\section{Memory Jogger}
Congress may use its Section 5 enforcement power to abrogate state sovereign immunity and authorize private damages and fee awards against unconsenting States (\href{https://supreme.justia.com/cases/federal/us/427/445/}{427 U.S. 445}).

\section{Detailed Case Facts}
Connecticut's retirement plan let women with 25 years of service retire on full pensions at 50, while men had to wait until 55. The 1972 Amendments to Title VII, enacted under Section 5 of the Fourteenth Amendment, extended Title VII to state employers and authorized backpay and attorney's fees. Male employees led by Garland Fitzpatrick sued state officials, including J. Frederick Bitzer, for retroactive benefits payable from the state treasury.

\section{Procedural History}
The district court found a Title VII violation and enjoined the plan but, under the Eleventh Amendment and \textit{Edelman v.\ Jordan}, denied backpay and fees. The Second Circuit affirmed on backpay but allowed fees. The Supreme Court granted cross-petitions.

\section{Judicial Votes}
9--0 in the judgment: Rehnquist wrote for seven Justices; Brennan and Stevens concurred in the judgment.

\section{Holding}
The Eleventh Amendment is limited by Section 5 of the Fourteenth Amendment. When Congress clearly authorizes suits under Section 5, federal courts may award retroactive backpay and attorney's fees against States.

\section{Analysis of Opinions}
Rehnquist reasoned that the Fourteenth Amendment, ratified after the Eleventh, pairs prohibitions directed at the States with an enforcement grant to Congress; under \textit{Ex parte Virginia}, the Civil War Amendments were ``limitations of the power of the States and enlargements of the power of Congress.'' Congress may therefore authorize suits impermissible in other contexts, and fees stand on the same footing. Brennan would have held that the Eleventh Amendment does not bar suits by a State's own citizens at all; Stevens would have relied on the Commerce Clause and a narrower reading of \textit{Edelman}. The Article I contrast foreshadowed \textit{Seminole Tribe}.

\section{Examples: Future Applications}
\textbf{Same-Side:} FMLA family-care damages against States, supported by a record of sex discrimination, are valid (\textit{Hibbs}). \textbf{Opposite-Side:} a Commerce Clause statute letting tribes sue States cannot abrogate immunity (\textit{Seminole Tribe}). \textbf{Fence-Sitter:} damages against state universities for disparate-impact sex discrimination, which \textit{Fitzpatrick} left open.

\section{Critique}
Originalists praise the last-in-time logic but note that the text never mentions damages suits and that the Court skipped whether Title VII tracked actual violations, later forcing \textit{City of Boerne}. Progressives argue the Eleventh Amendment never covered a State's own citizens and that, after \textit{Seminole Tribe}, \textit{Fitzpatrick} became a narrow bottleneck.

\section{Key Quotations}
\begin{quote}\sloppy
``But we think that the Eleventh Amendment, and the principle of state sovereignty which it embodies, see \textit{Hans v.\ Louisiana}, 134 U.S.\ 1 (1890), are necessarily limited by the enforcement provisions of \S~5 of the Fourteenth Amendment.''
\end{quote}
\noindent --- Justice Rehnquist, Opinion of the Court, 427 U.S.\ at 456.

\end{document}
